\section{Lazy-sampling QROs: The Compressed Oracle Technique}
  We will employ the compressed oracle technique, as introduced by
    Zhandry~\cite{C:Zhandry19};
    specifically, we employ an adapted version of the ``standard Fourier
    oracle'', the choices
    being standard/phase and computational/Fourier, depending on what
    oracle interface we want, and in which domain the database entries should
    be defined.\footnote{
      The standard Fourier oracle is not the version
      employed by Zhandry in his paper, and it is only defined in an appendix,
      but it is the version I find most closely
      matches the intuition behind the compressed oracle technique.
    }

  \subsection{Setup}
    Let us begin by specifying our notation. To start with, unless noted otherwise,
      we will assume that all random oracles have input length $n$ and
      output length $m$, and that $q$ is an upper bound on the total
      number of queries made to the oracle.

    \paragraph{The empty string.}
       Let $\emptystring$ to be the empty string, encoded
        as $1 \concat{}0^n$. We will interpret the string $0 \concat{}x$ as $x$.
        Strings of the form $1 \concat x$ for $x \neq 0^n$ will never
        appear, but for the sake of reference let us call them $\slashed{x}$.

      We define the $(1+n)$-bit controlled-not operation on $\emptystring$ to be
        $\eCXgate(\ket{x}_{\register[A]} \to \ket{\emptystring}_{\register[B]})
        = \ket{x}_{\register[A]}\ket{x}_{\register[B]}$.
        This may be implemented by performing the controlled not-operation as usual,
        and then flipping the leading bit of register $\register[B]$. Repeating the
        operation yields $\eCXgate(\ket{x}_{\register[A]} \to
        \ket{x}_{\register[B]}) =
        \ket{x}_{\register[A]}\ket{\emptystring}_{\register[B]}$ again, making
        $\eCXgate$ an involution.
        (Applying $\eCXgate$ to states $\ket{x}\ket{y}$ such that $x \neq y$
        will lead to the invalid state $\ket{x}\ket{\slashed{z}}$, where $z = x \xor
        y$, but we will never apply the gate to such cases.)

      Finally, observe that $x \cdot \emptystring = 0$ for any $x$, where $\cdot$
        denotes the bitwise inner product.

    \paragraph{Databases.}
      We are going to keep a database register $\database$, storing $q$
        list entries of the format $(b_\emptystring \concat x, \hat{y})$, meaning $\database$ is a
        register containing $q \cdot (1+n+m)$ qubits.
        The database starts out in the ``empty'' state, with each entry
        initialized in the pure state $\ket{\emptystring, 0^m}$.
        After several
        queries, the database will end up in a superposition over lists $\Dlist$, with
        each list being split in two parts: An ``active'' part of tupels
        $(x,\hat{y})$ such that $\hat{y} \neq 0^n$, sorted by increasing $x$, followed by
        ``fresh'' part of entries set to the initial value $(\emptystring, 0^m)$.
        More precisely, after $\ell$ queries, $\database$ will be in a
        superposition over basis states of the form
      \begin{align*}
        \ket{\Dlist} &= \ket{(x_1, \hat{y_1}), (x_2, \hat{y}_2) \ldots, (x_k,
        \hat{y}_k), (\emptystring, 0^m), \ldots, (\emptystring, 0^m)} \\
        &= \ket{x_1, \hat{y}_1}_1 \ket{x_2, \hat{y}_2}_2 \dots
        \ket{x_k, \hat{y}_k}_k \ket{\emptystring, 0^m}_{k+1} \dots
        \ket{\emptystring, 0^m}_q\, ,
      \end{align*}
      with unique and increasing input elements $x_1 < x_2 < \ldots < x_k$,
        non-zero output elementst $\hat{y}_i$, and $k \leq \ell$.
        We write $\Dlist[x] = \hat{y}$ if the list contains an element
        $(x,\hat{y})$, and $\Dlist[x] = \perp$ otherwise. 
        %This is well-defined,
        %since we are guaranteed that the list only houses at most one tuple with
        %the value $x$ at any time.

      The state of the database register will in general be highly entangled
        with the parties performing the queries; indeed, a list element will be
        non-empty \emph{only if} the register is entangled with the state of a
        querying party. Finally, the number of ``active'' versus ``fresh''
        entries may differ between branches of the superposition of $\database$,
        though their sum will always be $q$.\footnote{
          In the original formulation of the compressed oracle technique, the size
          of $\database$ was allowed to grow and shrink with each call.
          This leads to a superposition over databases where also the
          \emph{size} of the register is in superposition. In our formalization,
          registers are classically well-defined lists storing qubits (which are
          two-level physical systems), and
          while it is possible to resize registers by moving qubits around, we do
          not allow for this to happen in superposition. Therefore, we fix
          the size of $\database$ to be $q$ a priori,
          and simply accept that the simulation will fail if the query count
          ever exceeds $q$.
        }

    \paragraph{Database subregisters.}
      Next, let us define a naming convention for the subregisters of $\database$:
        Let $\database_\inputreg$ refer to the collection of qubits that
        store input values in the list (i.e.\ a subregister of size $q\cdot
        (1+n)$), and let $\database_\outputreg$ refer to the collection of
        qubits that stores output values (with size $q \cdot m$).

      Now, provided $\Dlist[x] \neq \emptystring$,
        let $\database(x)$ refer to the qubits that store the
        relevant tuple $(x, \hat{y})$. Note that even in branches where this
        holds, the exact qubit indices corresponding to
        $\database(x)$ will in general differ between branches, since the lists in the
        superposition may contain a varying number of entries that precede $x$,
        making $\database(x)$ only well-defined when specializing to single
        branches of $\database_\inputreg$ superposition.
        When speaking generally, we may use $\database_i$ to refer to the registers
        that will store the $i$'th tuple of any list.

      Combining, we can then refer to the subregister storing the first element of
        the tuple $(x,\hat y)$ as $\database_\inputreg(x)$, and the register
        storing the corresponding value $\hat{y}$ as
        $\database_\outputreg(x)$, (and likewise for $\database_{\inputreg,i}$,
        $\database_{\outputreg,i}$).

      %\hhnote{
      %  I \emph{think} \ldots but it may be the case (indeed, the \emph{point})
      %  that all empty elements are pure states, i.e.\ computational basis
      %  states, in the Hadamard domain. Still, though, it seems like we have no
      %  choice but to operate element-wise \ldots can you even recognize a
      %  pure basis state without measuring it? In that sense, the real
      %  ``empty'' element will be the inputs that remain undefined in
      %  \emph{all} lists in the database superposition. If the adversary
      %  forwards a uniform superposition as input, then there will be no such
      %  state: We will have a uniform superposition over $2^n$ lists, each of which
      %  contains the single element $(i, 0^n)$ at the point of returning.
      %  (Still not completely wrapped my head around how the second part gets
      %  updated away from $0^n$. Something something entanglement\ldots)
      %}

    \paragraph{The primal and the Hadamard domain.}
      You may have noticed the hats. These denote that the value in question is
        defined in the Hadamard-domain: A computational basis state
        $\ket{y}$ in the primal domain (left) will be the following 
        superposition in the Hadamard domain (right):
        \begin{align*}
          \ket{y}
          \quad &\applyleftright{H^{\otimes m}} \quad
          \sum_{\forall \hat{y} \in [2^m]} \frac{(-1)^{y \cdot
          \hat{y}}}{\sqrt{2^m}} \ket{\hat{y}}\, ,
          \intertext{
            and vice versa:
          }
          \sum_{\forall y \in [2^m]} \frac{(-1)^{\hat{y} \cdot y}}{\sqrt{2^m}} \ket{y}
          \quad &\applyleftright{H^{\otimes m}} \quad
          \ket{\hat{y}}\, .
          \intertext{
            In particular, we have that the all-zero state in the Hadamard domain
            corresponds to the uniform superposition in the primal domain:
          }
          \sum_{\forall y \in [2^m]} \frac{1}{\sqrt{2^m}} \ket{y}
          \quad &\applyleftright{H^{\otimes m}} \quad
          \ket{0^m}\, .
        \end{align*}

      \noindent
      As illustrated, one can move between the primal domain and the
        Hadamard domain by performing the ``Hadamard transformation'',
        which is just applying the $\Hgate$ gate to each qubit in the
        register individually.\footnote{
          In his paper~\cite{C:Zhandry19}, Zhandry varyingly refers to the Quantum Fourier
          Transformation and the Hadamard transformation, but he clarifies in
          App.\ A that he only ever uses the latter. (The Hadamard gate $\Hgate$ is the
          single-qubit Quantum Fourier Transformation.)
        } We will only move between domains on the output register,
        i.e.\ the $\outputreg$ register forwarded as part of the query.  
        The input registers $\inputreg$ and
        $\database_{\inputreg}$ will both be kept in the primal domain,
        and the output register $\database_{\outputreg}$ of the database 
        will be kept in the Hadamard domain.

  \subsection{Defining CROs}
    \paragraph{Updating the database.}
      Next, we define a unitary that acts on the combined state of the input
        register $\inputreg$ and the database $\database$,
        named $\decomp$ for ``Store/Sort''.\footnote{
          Zhandry calls this unitary $\unitary[FourDecomp]$, for ``Fourier
          Decompress''~\cite[App.~C]{C:Zhandry19}.
        } In essence, it is an operation that
        reads the input, and updates the database accordingly.
        We define it by specifying it on all basis states, verifying along the
        way that each step is reversible, and therefore
        unitary. Indeed, $\decomp$ will turn out to be an involution, meaning
        that applying it twice in a row will return the same state, (or in 
        other words, that $\decomp^2 = \idmatrix$).

      \begin{definition}[Unitary $\decomp$]
        For each computational basis state $\ket{x}\ket{\Dlist}$ in the
          support of the combined state of registers $\inputreg$ and $\database$:
        \begin{enumerate}
          \item Do one of the following, depending on the case: 
          \begin{enumerate}
            \item If $\Dlist[x] = \hat{y} \neq 0^m$, then $\RO(x)$ has already been
              set, and $\database(x)$ must furthermore be entangled with a
              querying party; do nothing.
            \item If $\Dlist[x] = \perp$, set
              the first empty element $(\emptystring, 0^m)$ in $\Dlist$ to
              $(x, 0^m)$ by applying $\eCXgate$ from $\inputreg$ to the
              corresponding subregister $\database_{\inputreg,i}$.
            \item If $\Dlist[x] = 0^m$, ``forget'' the element by applying 
              $\eCXgate$ from $\inputreg$ to the subregister
              of $\database_{\inputreg}(x)$,
              resulting in the entry $(\emptystring, 0^m)$. 
              %(Notice that this
              %operation is identical to the above one, but with the inverse outcome.)
            \item For all other basis states, which is to say for lists $\Dlist$ that 
              contain no empty entries $(\emptystring, 0^m)$ (which can
              only happen if the query count exceeds $q$), and for lists that
              contain invalid values $\slashed{x}$ for some $x \neq 0^n$ (which will
              never happen), act as the identity.
          \end{enumerate}
          \item Use controlled $\SWAPgate$ gates to sort $\Dlist$ by increasing $x$.
        \end{enumerate}

        \noindent
        The If-statements may be implemented as mult-bit controlled
        operations. Notice that the cases are mutually exclusive as well as
        exhaustive, and that they define an involution. Therefore, $\decomp$ is
        unitary.
      \end{definition}

    \paragraph{Entangling the query.}
      With the updating procedure defined, all that remains is to 
        define how an output is provided. 

      In essence, the idea is to
        move the output register $\outputreg$ to the Hadamard
        domain, entangle it with the subregisters $\database_{\outputreg}(x)$
        for each $x$ in the input superposition,
        and then return $\outputreg$ to the primal domain, before shipping
        $\inputreg$ and $\outputreg$ back to the querying party.
        The entanglement, which is hidden from view of whoever made the query,
        will make it look like a classical probability distribution has been
        induced over the possible states of $\outputreg$---just as if $\RO$ had
        been sampled from the set of all functions.

      Let us define the ``Entangling'' unitary $\entangle$ next.

      \begin{definition}
        Let $\entangle$ be the unitary that acts on basis states
          $\ket{x}\ket{y}\ket{\Dlist}$ of registers
          $\inputreg$, $\outputreg$, $\database$
          in the following way:
          \[
            \ket{x}\ket{\hat y}\ket{\ldots, (x, z), \ldots} \apply{\entangle}
            \ket{x}\ket{\hat y}\ket{\ldots, (x, \hat y \xor z), \ldots}\, ,
          \]
          and as the identity whenever $\Dlist[x] = \perp$.\footnote{
            We will ensure that this never occurs.
          } The cases are mutually exclusive and exhaustive, and define
          an involution. Therefore, $\entangle$ is unitary.
      \end{definition}

    \paragraph{The oracle procedure.}
      With $\decomp$ and $\entangle$ defined, we are finally ready to define
        the action performed by our compressed random oracle upon receiving a query.

      \begin{definition}[Compressed random oracle]
        \label{def:compressed-ros}
        Given registers $\inputreg,
          \outputreg, \database$, the \emph{compressed random oracle} 
          $\RO: \bin^n \to \bin^m$ performs the following operations:
        \begin{enumerate}
          \item $\applyU{\decomp}{\inputreg, \database}$
          \item $\applyU{\Hgate^{\otimes m}}{\outputreg}$
          \item $\applyU{\entangle}{\inputreg, \outputreg, \database}$
          \item $\applyU{\Hgate^{\otimes m}}{\outputreg}$
          \item $\applyU{\decomp}{\inputreg, \database}$
        \end{enumerate}
      \end{definition}

      \noindent
      Having defined this procedure, Zhandry then proved that it provides a perfect
        simulation of a quantum-accessible random
        oracle~\cite[Lemma~4]{C:Zhandry19}.

      \begin{theorem}[Simulation is perfect]\label{thm:perfect}
        Let $\RO$ be a quantum-accessible random oracle, and let $\RO[G]$ be a
        compressed random oracle. Then we have for all distinguishers $\distinguisher$,
        \[
          \prob{1 \sample \distinguisher^{\RO}} = \prob{1 \sample
          \distinguisher^{\RO[G]}}\, .
        \]
      \end{theorem}

  \subsection{Playing with CROs}
    What is the effect of the operations listed in \autoref{def:compressed-ros}? Let's explore!

    \paragraph{Classical queries.}
      For our first sanity check, consider the result of performing the
        first three steps on a classical query, i.e.\ with $\inputreg$
        and $\outputreg$ put in a computational basis state
        $\ket{x}\ket{y}$, to the initially-empty database
        $\ket{\Dlist_\emptystring}$:
        \begin{enumerate}
          \item $\ket{x}\ket{y}\ket{\Dlist_\emptystring} \coloneqq
            \ket{x}\ket{y}\ket{(\emptystring, 0^m), \ldots, (\emptystring,
            0^m)}%$ \\ 
            \quad \apply{\decomp} \quad \ket{x}\ket{y} \ket{(x, 0^m), (\emptystring, 0^m),
            \ldots, (\emptystring, 0^m)} \coloneqq \ket{x}\ket{y}\ket{\Dlist_x}\, ,$
          \item $\ket{x}\ket{y}\ket{\Dlist_x}%$ \\
            \quad \apply{H^{\otimes m}} \quad \ket{x}
            \left(\sum_{\forall \hat{y} \in [2^m]} \frac{(-1)^{y \cdot \hat{y}}}{\sqrt{2^m}}
            \ket{\hat{y}}\right) \ket{\Dlist_x}\, ,$
          \item $\ket{x} \left(\sum_{\forall \hat{y} \in [2^m]} \frac{(-1)^{y
            \cdot \hat{y}}}{\sqrt{2^m}} \ket{\hat{y}}\right) 
            \ket{(x, 0^m), (\emptystring, 0^m), \ldots}$ \\
            $\apply{\entangle} \quad 
            \ket{x}\left(\sum_{\forall \hat{y} \in [2^m]}
              \frac{(-1)^{y \cdot \hat{y}}}{\sqrt{2^m}} \ket{\hat{y}} \ket{(x, \hat{y}),
              (\emptystring, 0^m), \dots} \right)
            \coloneqq \sum_{\forall \hat{y} \in [2^m]}
              \frac{(-1)^{y \cdot \hat{y}}}{\sqrt{2^m}}
              \ket{x}\ket{\hat{y}}\ket{\Dlist_{x,\hat{y}}}\, .$
        \end{enumerate}

      \noindent
      At this point, $\outputreg$ is maximally entangled with the
        subregister of $\database$ housing $\Dlist[x]$: 
        We can see from the above that their combined
        state now has the form (ignoring amplitudes) 
        $\ket{\hat{y}_1}\ket{\hat{y}_1} +
        \ket{\hat{y}_2}\ket{\hat{y}_2} + \ldots +
        \ket{\hat{y}_{2^m}}\ket{\hat{y}_{2^m}}$. This may be a large superposition,
        but the outcome of observing the two registers is guaranteed to be
        equal, which is exactly what it means to be maximally entangled (in
        general, that observing one register fully determines the state of the other).
        Performing the Hadamard transformation on $\outputreg$ thereafter does nothing to change
        this\footnote{
          This is because entanglement is a \emph{basis-independent} feature of
          quantum states; in fact, it fulfills the definition of a \emph{quantum resource},
          i.e.\ something that can be spent to perform tasks~\url{https://arxiv.org/abs/1806.06107}.
        }, but rather has the effect of moving their combined state from one maximally
        entangled state to another.%by simply shuffling phases around. 

      Finally, the second application of
        $\decomp$ does nothing, since the list is already sorted and there is
        nothing to remove. That is, with one exception: In the single branch of the
        wavefunction where $\hat{y} = 0^m$, $x$ will be replaced with
        $\emptystring$, which has the effect of ``implicitly'' defining the
        value of $\hat{y}$ in this branch. 

      At a glance, this may appear to be a
        problem, since if the querying party then by chance observes exactly
        this branch, and then queries the compressed oracle again on the same
        input, the database would not ``know'' that it was supposed give the same
        value back, leading to an inconsistency. However, \autoref{thm:perfect}
        tells us that simulation is perfect, so \emph{something} will have
        to ensure consistency in this case. The answer will be
        given in terms of the $m=1$ case in \autoref{sec:extracting}, though I
        urge you to work it out for yourself and seeing what happens before
        checking your answer.

      And so, accepting this conundrum for now, we return to our maximally entangled
        state. The entanglement is hidden from view of the querying party, and
        hiding away entanglement in this manner has the effect of ``mixing''
        the state of $\outputreg$: From their point of view, hidden
        entanglement is indistinguishable from a classical probability
        distribution over quantum states\footnote{
          In fact, this is the observation that lead Hugh Everett to
          develop his many-worlds interpretation of quantum mechanics,
          since the randomness induced by wavefunction collapse can just
          as well be explained in terms of you entangling \emph{yourself} with
          the quantum state you are observing. And why not? After all, you,
          too, are made of quantum stuff!
        }; and maximal entanglement yields
        maximally mixed states, which corresponds to a uniform probability
        distribution over basis states.
        
      Thus, the querying party has the experience of having the
        $\outputreg$ register placed in a uniformly random computational
        basis state---exactly as if $\RO(x)$ had been sampled at random
        and xored into the register!

      %\hhnote{
      %  What is the effect of the $0^m$ part of the superposition being removed
      %  by $\decomp$? (I think it is not ``removed'' per se, but rather
      %  implicitly defined, since we don't have to keep track of $x$ for such
      %  cases. The output register is still \emph{there}, after all.)
      %}

    \paragraph{Uniform superpositions.}
      Next, what if $\outputreg$ is set to be in the
        uniform superposition, $\sum_y \sqrt{2^{-m}}\ket{y}$ (with the input
        value still set to the computational basis state $\ket{x}$)? Well, then
        \emph{nothing} happens: After being moved to the Hadamard
        domain, the $\outputreg$ register will be in the state
        $\ket{0^m}$, and xoring into the database has no effect: The second
        Hadamard transformation will take it back to the uniform superposition,
        and then return it as-is. Note that this also matches the querying party's
        expectations: Since the xor of any fixed string defines a permutation, xoring 
        into the uniform superposition should yield the uniform superposition,
        regardless of the value $\RO(x)$ that is being xored.
        This means that, when interacting with a real random oracle, $\RO(x)$
        will remain perfectly hidden from their view. This is matched in the compressed
        oracle by the second application of $\decomp$ removing the point $x$
        from the list and replacing it with $\emptystring$, ensuring no
        entanglement (which would otherwise have been a noticeable
        inconsistency, since the output register would no longer be in a
        uniform superposition, but rather in some mixed state).

      What about other superpositions, for example $\sum_y \nicefrac{(-1)^{y
        \cdot z}}{\sqrt{2^{-m}}} \ket{y}$ for some $z \in \bin^m$? Well, then the querying party
        \emph{would} expect to be able to learn something about the output, by
        observing how the phases shift around. The database will likewise be
        updated in this case: In our example, moving $\outputreg$ to the Hadamard domain
        will put it in the computational basis state $\ket{z}$, which
        will $\entangle$ will then store into the database as $(x,z)$. Now we
        are back to the case of maximal entanglement, as can be seen from the
        combined state of $\outputreg$ and $\database_\outputreg(x)$ being
        $\ket{z}\ket{z}$.

      Of course, in between the two extremes of no information and maximal
        entanglement lies the world of \emph{some} entanglement, as
        results from setting $\outputreg$ to some non-uniform superposition.
        Additionally, $\outputreg$ may itself be arbitrarily entangled with
        other subsystems belonging to the querying party, in which case it will
        be a mixed state from the point of view of the compressed oracle, and
        so their combined state will, after entangling, also be a mixed state.

    \paragraph{Superpositions over input values.}
      Finally, what about superpositions over $x$? The easiest way to think
        about this is to adopt an Everettian worldview, and imagine the
        superposition state of $\inputreg$
        describing a superposition over several possible universes. Then, one can
        consider each universe separately: The branch of the wavefunction where
        $\inputreg$ is in basis state $\ket{x}$ is the branch of the
        wavefunction where the querying party (who will in general be entangled
        with $\inputreg$) decided to query the random oracle on the value $x$. 
        And since these registers always remain in the primal domain, we don't have
        to grapple with the metaphysical implications of ``changing the basis'' on a
        collection of universes.

      Thus, the result of such a query again depends on
        the state of $\outputreg$: If $\outputreg$ is in
        the uniform superposition, independently of $\inputreg$,
        then it is in the uniform superposition in every universe, and
        nothing is ever stored and nothing is ever learned. If
        $\outputreg$ is set in a computational basis state, such as
        $\ket{0^m}$, then the result is a maximally entangled state,
        appearing as a classically uniform probability distribution over basis
        states in all universes. Finally, if $\inputreg$ and $\outputreg$ are
        maximally entangled in some state $\ket{x_1}\ket{y_1} + \ket{x_2}\ket{y_2} +
        \ldots$, for unique values of $x_i$ and $y_i$, then
        inside each universe we are performing a purely classical query,
        and we are back to the situation that we first considered.

    \paragraph{Repeated queries.}
      Say the oracle started out in the empty state
        $\ket{\database_\emptystring}$, and was then queried on the classical
        input state $\ket{x}\ket{y}$.
        Then we found that the result is the combined state
      \[
        \sum_{\forall \hat{y} \in [2^m]\backslash \{0^m\}} \frac{(-1)^{y \cdot \hat{y}}}{\sqrt{2^m}}
        \ket{x} \left( \Hgate^{\otimes m} \ket{\hat{y}} \right) \ket{\Dlist_{x,\hat{y}}}
        + \frac{1}{\sqrt{2^m}} \ket{x} \left( \Hgate^{\otimes m} \ket{0^m} \right) \ket{\Dlist_\emptystring}\, . 
      \]

      \noindent
      Now, what happens if the adversary queries the oracle on the registers
        $\inputreg$ and $\outputreg$ a second time, without changing their
        in any way in the meantime? Well, first, the Hadamard transformation on
        $\outputreg$, being an involution, cancels out:
      \[
        \sum_{\forall \hat{y} \in [2^m]\backslash \{0^m\}} \frac{(-1)^{y \cdot \hat{y}}}{\sqrt{2^m}}
        \ket{x} \ket{\hat{y}} \ket{\Dlist_{x,\hat{y}}}
        + \frac{1}{\sqrt{2^m}} \ket{x} \ket{0^m} \ket{\Dlist_\emptystring}\, . 
      \]

      \noindent
      Next, for each basis state $\ket{x} \ket{\hat{y}}
        \ket{\Dlist_{x,\hat{y}}}$ such that $\hat y \neq 0^m$, we will have
        that $\Dlist[x] \neq \perp$, and so $\decomp$ will do nothing,
        while for $\hat y = 0^m$ it will insert the element $(x, 0^m)$ into
        $\Dlist$. This gives us the state:
      \[
        \sum_{\forall \hat{y} \in [2^m]} \frac{(-1)^{y \cdot \hat{y}}}{\sqrt{2^m}}
          \ket{x} \ket{\hat{y}} \ket{\Dlist_{x,\hat{y}}}\, .
      \]

      \noindent
      Then, $\entangle$ will have the following effect on the basis states:
        \[
          \ket{x}\ket{\hat{y}}\ket{(x, \hat{y}), \ldots} \apply{\entangle}
          \ket{x}\ket{y}\ket{(x, \hat{y} \xor \hat{y}), \ldots}
          = \ket{x}\ket{y}\ket{(x, 0^m), \ldots}\, ,
        \]

      \noindent
      and with the second application of $\decomp$ removing all entries with
        $\hat{y} = 0^m$, we will end up with the unentangled state
      \[
        \sum_{\forall \hat{y} \in [2^m]} \frac{(-1)^{y \cdot \hat{y}}}{\sqrt{2^m}}
          \ket{x} \ket{\hat{y}} \ket{\Dlist_\emptystring}\, .
      \]

      \noindent
      Transforming $\outputreg$ back to the primal domain then yields
      \[
          \ket{x} \ket{y} \ket{\Dlist_\emptystring}\, ,
      \]
      which is exactly where we started before the first query!

      We see that the entanglement has been removed,
        and the value of $x$ has been completely removed from the oracle's view.
        What's more, since this ``forgetting'' behaviour is shown to hold for
        any computational basis state, it must hold for \emph{every} quantum
        state that is forwarded to the oracle twice in a row, making the
        compressed oracle \emph{itself} an involution!
        Worse, not only does the oracle seem to ``forget'' that it was ever
        queried, but so does the querying party also ``forget'' the value of
        $\RO(x)$!

      It may seem counter-intuitive, but this ``forgetting'' is in fact exactly
        what we need to happen (and was indeed a major roadblock towards
        lazy-sampling quantum random oracles prior to the development of this
        technique). To see why, consider the situation for
        ``monolithic'' random oracles, which are also involutions: 
        Here, the classical query $\ket{x}\ket{y}$ returns the state $\ket{x}\ket{y \xor
        \RO(x)}$. If the query is then repeated with the same registers
        without altering their state, the result will be $\ket{x}\ket{y \xor \RO(x)\xor \RO(x)} =
        \ket{x}\ket{y}$. The value of $\RO(x)$ has been forgotten!

      ``But'', I hear you say, ``what if I store the value of $\RO(x)$ before
        I make the second query? Surely I won't be forced to automatically
        `forget' what the value was, simply because I query the oracle a second
        time!'' And you'd be right---but notice that we are now talking about a
        different quantum state, namely the state
        $\ket{x}_\inputreg \ket{y\xor \RO(x)}_\outputreg
        \ket{\RO(x)}_{\register[A]}$, where the $\register[A]$ register is
        some storage register local to you. By storing $\RO(x)$ in this
        register (by first preparing it in the state $\ket{y}$ and
        then performing the $m$-bit $\CXgate$ gate from $\outputreg$ to
        $\register[A]$), you \emph{entangled} the two registers. And so,
        while the result of a second query to the monolithic oracle will indeed be
        the state $\ket{x}\ket{y}\ket{\RO(x)}$, the effect on
        the compressed oracle will be completely different to the
        ``forgetting'' behaviour outlined above. This can be seen from the fact
        that the Hadamard transformation, which previously operated on the
        pure, unentangled state $\ket{y}$, will now take the $\outputreg$ register
        to some other state than the one it was in in during the first query, and the
        cancellation of the $\hat{y}$ values in the second query won't happen.

    \subsection{Extracting Queries}\label{sec:extracting}
      From the point of view of the
        querying party, should they observe this value, it will be the same as
        seeing $\RO(x) = 0^m$; this will still happen with probability $2^{-m}$.
        However, since $x$ was removed from the database, making the same query
        again will lead to a completely different result with overwhelming
        probability.

      \hhnote{
        But the database won't know this, and querying again on the same point
        will lead to a different outcome with overwhelming probability. Is this
        a problem? Sure \emph{seems} like a problem! But $\hat{y} = 0^m$ is a
        uniform superposition in the primal domain \ldots but this is a
        maximally entangled state \ldots I am confused. \\
        I think that the answer is that in this formalization, there is
        \emph{one} output value which will lead to inconsistent answers upon
        repeated queries post-measurement, but that its exact value (input and
        output alike) will be perfectly unpredictable.
      }
